\documentclass[10pt,a4paper]{amsart}
\usepackage[margin=19mm]{geometry}
\usepackage{amsmath,amssymb,mathtools}
\usepackage[scr=rsfs,cal=euler]{mathalfa}
\usepackage{xfrac}
\usepackage[unicode,hidelinks,bookmarks=false]{hyperref}
\newcommand{\ie}{i.e.\ }
\newcommand{\cf}{cf.\ }
\newcommand{\resp}{respectively\ }
\newcommand{\Subsection}{\subsection}
\providecommand{\nobreakdash}{\nobreak\hbox{-}}
\setlength{\emergencystretch}{2em}
\multlinegap0pt
\usepackage{amssymb}
\usepackage{bm}
\newtheorem{theorem}{Theorem}
\newtheorem{lemma}{Lemma}
\newcommand{\abs}[1]{\left\lvert #1\right\rvert}
\newcommand{\bbz}{\mathbb{Z}}
\newcommand{\brac}[1]{\left( #1\right)}
\newcommand{\ogcd}[2]{\gcd(#1,#2)}
\title{Odd moments and adding fractions}
\author{Thomas F. Bloom, Vivian Kuperberg}
\date{Source: arXiv:2312.09021v2 (CC BY 4.0)}
\begin{document}
\maketitle

\noindent\emph{Source and licence.} Odd moments and adding fractions. Version arXiv:2312.09021v2 (CC BY 4.0), distributed by arXiv under the Creative Commons Attribution 4.0 International licence, http://creativecommons.org/licenses/by/4.0/, as recorded in the arXiv licence metadata for this exact version and shown on the abstract page https://arxiv.org/abs/2312.09021v2. Full licence notice: \texttt{licenses/arxiv-2312.09021-CC-BY-4.0.txt}.\par\smallskip

\noindent\textit{Editorial scope.} Counted unit: the article's lemma lem:bounds-on-s-i-sums-with-alpha-shifts (source0.tex lines 479--484). The article's complete original proof of it is delivered with it (source0.tex lines 485--512, one \texttt{proof} environment of 28 lines). No mathematics is written, repaired or re-derived: the statement and the proof are the article's own lines, in the article's own notation, and no retained line is substituted or re-flowed.  Retained uncounted as the source-local context the proof needs: lines 159--166.  Nothing was omitted from any retained span, so the package carries no omission record.  No retained line contains a citation, so the wrapper carries no bibliography and no bibliography entry.  No retained line contains a figure environment, an image-inclusion command or drawing code, so no image is delivered and the package declares no asset dependency.  Editorial formatting only, in addition: an AMS \texttt{amsart} preamble that restates the article's own declarations and macros used by the retained lines, namely the environments \texttt{theorem}, \texttt{lemma} on separate counters, together with the standard packages those lines need.  The retained environments print in this document as Theorem 1, Lemma 1 in order of appearance, whereas the article prints the same items with the numbering of its own full document.  The complete native source of the article as distributed in the arXiv e-print for this exact version is bundled unchanged as \texttt{sources/151506/source0.tex}.
\begin{theorem}\label{th-oddsum}
Let $k\geq 1$ be an integer. Let $Q_1,\ldots,Q_{2k+1}\geq 1$ and $A_1,\ldots,A_{2k+1}\subseteq [-1,1]$ be intervals of lengths $\delta_1,\ldots,\delta_{2k+1}$ respectively, where $\delta_i\geq 1/Q_i$ for $1\leq i\leq 2k+1$. 

There exists some $X\subseteq \{1,\ldots,2k+1\}$ with $\abs{X}=k+1$ such that the number of solutions to
\[\frac{a_1}{q_1}+\cdots+\frac{a_{2k+1}}{q_{2k+1}}\in \bbz\quad\textrm{ where }\quad\ogcd{a_i}{q_i}=1\textrm{ for }1\leq i\leq 2k+1,\] 
with $a_i/q_i\in A_i$ and $q_i\leq Q_i$ for $1\leq i\leq 2k+1$, is
\[\ll (\log(Q_1\cdots Q_{2k+1}))^{O(1)}\brac{\prod_{i\in X}\delta_i}Q_1\cdots Q_{2k+1}.\]
\end{theorem}

\begin{lemma}\label{lem:bounds-on-s-i-sums-with-alpha-shifts}
Let $h \ge 1$ be an integer, let $k \ge 3$ be an odd integer, and fix any rational $\alpha_1, \dots, \alpha_k \in [0,1)$. Suppose $2\leq y\leq h$. For each $i$, let $\ell_i = 1$ if $\alpha_i = 0$ and let $\ell_i = h$ otherwise. Then
\begin{equation*}
    \sum_{\substack{s_1, \dots, s_k \\ \ell_i < s_i < y}} \frac{\mu(s_i)^2}{\phi(s_i)} \sum_{\substack{a_1, \dots, a_k \\ 1 \le a_i \le s_i \\ \ogcd{a_i}{s_i} = 1 \\ \sum_i a_i/s_i \in \mathbb Z}} \prod_{i=1}^k F\left(\frac{a_i}{s_i} + \alpha_i\right) \ll (\log y)^{O(1)}h^{\tfrac{k-1}{2}}.
\end{equation*}
\end{lemma}

\begin{proof}
The left-hand side can be bounded above by
\begin{equation}\label{eq-lem1}
\ll (\log y)^k \sum_{1 < s_1, \dots, s_k < y} (s_1 \cdots s_k)^{-1} \sum_{\substack{a_1, \dots, a_k \\ 1 \le a_i \le s_i \\ \ogcd{a_i}{s_i} = 1 \\ \sum_i a_i/s_i \in \mathbb Z}} \prod_{i=1}^k F\left(\frac{a_i}{s_i} + \alpha_i\right),
\end{equation}
since $\frac{s}{\phi(s)} \ll \log y$ for any $s \le y$. For $1 \le a_i \le s_i$, we have
\begin{equation*}
F\left(\frac{a_i}{s_i} +\alpha_i \right) = \min\left\{h, \left\|\frac{a_i}{s_i}+\alpha_i\right\|^{-1}\right\},
\end{equation*}
so that $\frac 1{s_i} F\left(\frac{a_i}{s_i}+\alpha_i\right) \le 1/n_i$, where
\begin{equation*}
n_i = \mathrm{max}\left\{\frac{s_i}{h}, \lfloor |\tilde{a_i} + \alpha_i s_i|\rfloor \right\},
\end{equation*}
and $\tilde{a_i}$ is the representative of $a_i$ mod $s_i$ which minimises $|\tilde{a_i} + \alpha_i s_i|$ (so that $|\tilde{a_i} + \alpha_i s_i|\leq s_i/2$). 

The possible values of $n_i$ are thus the set $\left\{\frac{s_i}{h}\right\} \cup \left(\mathbb Z \cap \left[\frac{s_i}{h},\frac{s_i}{2}\right]\right)$. 
Note that if $s_i/h < 1$, then $1 < s_i < h$, so $\alpha_i = 0$ by assumption. In this case there is no contribution when $n_i = s_i/h$, since that implies $|\tilde{a_i}| \le s_i/h < 1$, and thus $\tilde{a_i} = 0$, which contradicts the requirement that $\ogcd{a_i}{s_i} = 1$. In particular, this means that $n_i \ge 1$ always. 


Notice that since $a_i\leq s_i$ and $\lvert \tilde{a_i}+\alpha_is_i\rvert \leq n_i+1$, if we assume that $s_i\in [S_i,2S_i]$ say, there must exist some interval $I_i$, independent of $a_i$ and $s_i$, of length $O(n_i/S_i)$, which contains $a_i/s_i$. It follows from the discussion thus far that \eqref{eq-lem1} is bounded above by 
\[(\log y)^{O(1)}\sum_{\substack{n_1,\ldots,n_k\\ 1\leq n_i\leq y}}\prod_{i=1}^kn_i^{-1}\sup_{\substack{S_1,\ldots,S_k\\ 1\leq S_i\leq 2hn_i}} \Bigl(\sum_{\substack{s_1, \dots, s_k \\s_i\in [S_i,2S_i]}}\sum_{\substack{a_1, \dots, a_k \\ a_i/s_i \in I_i \\ \ogcd{a_i}{s_i} = 1 \\ \sum_i a_i/s_i \in \mathbb Z}}1\Bigr). \]
Noting that the term inside the brackets depends only on the size of each $n_i$, we can dyadically pigeonhole this sum according to the size of $n_i$ (at an additional cost of $(\log y)^{O(1)}$, and so it suffices to show that 

\[ \sup_{\substack{N_1,\ldots,N_k\\ 1\leq N_i\leq 2y}}\prod_{i=1}^kN_i^{-1}\sup_{\substack{S_1,\ldots,S_k\\ 1\leq S_i\leq 4hN_i}} \Bigl(\sum_{\substack{s_1, \dots, s_k \\s_i\in [S_i,2S_i]}}\sum_{\substack{a_1, \dots, a_k \\ a_i/s_i \in I_i \\ \ogcd{a_i}{s_i} = 1 \\ \sum_i a_i/s_i \in \mathbb Z}}1\Bigr)\ll (\log y)^{O(1)}h^{\lfloor k/2\rfloor},\]
where each $I_i$ is an interval of length $O(N_i/S_i)$ centred at $\alpha_i$. By Theorem~\ref{th-oddsum} the term inside the brackets is, for some $X\subseteq \{1,\ldots,k\}$ of size $\abs{X}=\lceil k/2\rceil$, 
\[\ll (\log y)^{O(1)}\prod_{i\in X}N_i\prod_{j\not\in X}S_j\ll (\log y)^{O(1)}h^{\lfloor k/2\rfloor} N_1\cdots N_k,\]
since each $S_i\ll hN_i$. This concludes the proof.
\end{proof}

\end{document}
